%% chemfrag-core.code.tex -- expandable engine of the chemfrag package
%%
%% Copyright (C) 2026 Vaibhav Blayer
%%
%% This work may be distributed and/or modified under the conditions of the
%% LaTeX Project Public License, either version 1.3c of this license or (at
%% your option) any later version. The latest version of this license is in
%%   https://www.latex-project.org/lppl.txt
%%
%% This work has the LPPL maintenance status `maintained'.
%% The Current Maintainer of this work is Vaibhav Blayer.
%%
%% ---------------------------------------------------------------------------
%% How it works
%% ---------------------------------------------------------------------------
%% chemfig expands a submolecule "!{name}" (or "!\macro") twice when it reads
%% it.  Every chemfrag fragment is a submolecule whose code is
%%      !\chemfragring{<ring>}{<keys>}   or   !\chemfragchain{<keys>}
%% Those two macros are *fully expandable*: their second expansion step
%% (\exp:w = \romannumeral) runs the whole engine below inside one \expanded
%% and leaves nothing but plain chemfig code.  chemfig then parses that code
%% as if the user had typed it, so fragments and hand-written chemfig mix
%% freely.
%%
%% Two rules follow from that design:
%%   * everything in this file is expandable;
%%   * chemfig reads its syntax characters ( [ ] : , - = ~ ( ) * @ ) with
%%     catcode 12, so every piece of syntax we emit is produced with
%%     \tl_to_str:n or a \c_..._str constant; user material (atoms,
%%     substituents) is passed through untouched with \exp_not:n.
%% ---------------------------------------------------------------------------

%% Defaults set with \chemfragset (searched after the user's keys).
\tl_new:N \l__chemfrag_defaults_tl

%% ===========================================================================
%% 1. Key lookup (expandable)
%% ===========================================================================
%% Key lists are ordinary "key=value" lists.  The first occurrence of a key
%% wins, so a fragment searches  <user keys> , <\chemfragset keys> , <data>.
%%
%% \__chemfrag_kv_all:nn {<list>} {<key>} leaves {<v1>}{<v2>}... for every
%% occurrence.  \keyval_parse:nnn wraps its result in \exp_not:n, so the
%% inner \expanded only *returns* the calls; they run in the surrounding
%% e-expansion.  This function must therefore always be used inside an
%% e-type expansion (\exp_args:Ne, \expanded, ...).
\cs_new:Npn \__chemfrag_kv_all:nn #1#2
  {
    \expanded
      {
        \keyval_parse:nnn
          { \__chemfrag_kv_flag:nn {#2} }
          { \__chemfrag_kv_pair:nnn {#2} }
          {#1}
      }
  }
\cs_new:Npn \__chemfrag_kv_flag:nn #1#2
  { \str_if_eq:nnT {#1} {#2} { { true } } }
\cs_new:Npn \__chemfrag_kv_pair:nnn #1#2#3
  { \str_if_eq:nnT {#1} {#2} { \exp_not:n { {#3} } } }

%% \__chemfrag_get:nnn {<list>} {<key>} {<default>}
%% Value of the first occurrence of <key>, protected against further
%% expansion; <default> when the key is absent.
\cs_new:Npn \__chemfrag_get:nnn #1#2#3
  { \exp_args:Ne \__chemfrag_get_aux:nn { \__chemfrag_kv_all:nn {#1} {#2} } {#3} }
\cs_new:Npn \__chemfrag_get_aux:nn #1#2
  {
    \tl_if_empty:nTF {#1}
      { \exp_not:n {#2} }
      { \exp_not:o { \use_i_delimit_by_q_stop:nw #1 \q_stop } }
  }

%% Same, but only "key=value" entries count (a bare "5" in !{chain}{5}
%% is the length, not a substituent on atom 5).
\cs_new:Npn \__chemfrag_kv_pairs:nn #1#2
  {
    \expanded
      {
        \keyval_parse:nnn
          { \use_none:n }
          { \__chemfrag_kv_pair:nnn {#2} }
          {#1}
      }
  }
\cs_new:Npn \__chemfrag_value:nnn #1#2#3
  { \exp_args:Ne \__chemfrag_get_aux:nn { \__chemfrag_kv_pairs:nn {#1} {#2} } {#3} }

%% Integer value of a key.
\cs_new:Npn \__chemfrag_int:nnn #1#2#3
  { \int_eval:n { \exp_args:Ne \use:n { \__chemfrag_get:nnn {#1} {#2} {#3} } } }

%% A flag is "on" when present with no value, or with any value but false.
\prg_new_conditional:Npnn \__chemfrag_if_on:nn #1#2 { T , F , TF }
  {
    \exp_args:Ne \str_if_eq:nnTF
      { \__chemfrag_get:nnn {#1} {#2} { false } } { false }
      { \prg_return_false: } { \prg_return_true: }
  }

%% Is <key> a bare integer (as in "!{chain}{5}")?  Leaves {<int>} if so.
\cs_new:Npn \__chemfrag_bare_int:n #1
  {
    \expanded
      { \keyval_parse:nnn { \__chemfrag_bare_int_aux:n } { \use_none:nn } {#1} }
  }
\cs_new:Npn \__chemfrag_bare_int_aux:n #1
  {
    \str_if_eq:eeT {#1}
      { \str_map_function:nN {#1} \__chemfrag_keep_digit:N }
      { {#1} }
  }
\cs_new:Npn \__chemfrag_keep_digit:N #1
  {
    \int_compare:nNnT { `#1 } > { 47 }
      { \int_compare:nNnT { `#1 } < { 58 } {#1} }
  }

%% Is the current key list member "<n>" present in the list stored under
%% <key>?  Used for  double={2,4}  and friends.
\prg_new_conditional:Npnn \__chemfrag_if_listed:nnn #1#2#3 { T , F , TF }
  {
    \exp_args:Ne \tl_if_empty:nTF
      {
        \exp_args:Ne \__chemfrag_kv_all:nn
          { \__chemfrag_get:nnn {#1} {#2} { } } {#3}
      }
      { \prg_return_false: } { \prg_return_true: }
  }

%% ===========================================================================
%% 2. Context coming from chemfig
%% ===========================================================================
%% A fragment that starts a molecule has no incoming bond.  chemfig keeps
%% track of what it did last in \CF_lastaction (0 = nothing yet).  If a
%% future chemfig drops that macro we simply assume an incoming bond.
\prg_new_conditional:Npnn \__chemfrag_if_standalone: { T , F , TF }
  {
    \cs_if_exist:NTF \CF_lastaction
      {
        \int_compare:nNnTF { \CF_lastaction } = { 0 }
          { \prg_return_true: } { \prg_return_false: }
      }
      { \prg_return_false: }
  }

%% ===========================================================================
%% 3. Emitting chemfig syntax
%% ===========================================================================
\cs_new:Npn \__chemfrag_hidden: { \exp_not:n { draw=none } }

%% \__chemfrag_bond:nnnn {<bond char>} {<relative angle>} {<length>} {<tikz>}
\cs_new:Npn \__chemfrag_bond:nnnn #1#2#3#4
  {
    #1 [ \c_colon_str \c_colon_str #2
    \tl_if_blank:nF { #3 #4 } { , #3 , , , #4 }
    ]
  }

%% Invisible bond, used to move to the centre of a ring and back out.
\cs_new:Npn \__chemfrag_jump:nn #1#2
  { \__chemfrag_bond:nnnn { - } {#1} {#2} { \__chemfrag_hidden: } }

%% Atom name for \chemmove:  @{<prefix><n>}
\cs_new:Npn \__chemfrag_name:nn #1#2
  { \exp_args:Ne \__chemfrag_name_aux:nn { \__chemfrag_get:nnn {#1} { name } { } } {#2} }
\cs_new:Npn \__chemfrag_name_aux:nn #1#2
  { \tl_if_blank:nF {#1} { \c_atsign_str { \exp_not:n {#1} #2 } } }

%% Content of atom <n>:  atoms={<n>=<content>, ...}
\cs_new:Npn \__chemfrag_atom:nn #1#2
  { \exp_args:Ne \__chemfrag_value:nnn { \__chemfrag_value:nnn {#1} { atoms } { } } {#2} { } }

%% Invisible copy of atom <n>, so that a bond leaving a heteroatom after an
%% exit jump starts at the edge of the letter, not at its centre.
\cs_new:Npn \__chemfrag_phantom:nn #1#2
  { \exp_args:Ne \__chemfrag_phantom_aux:n { \__chemfrag_atom:nn {#1} {#2} } }
\cs_new:Npn \__chemfrag_phantom_aux:n #1
  { \tl_if_blank:nF {#1} { \exp_not:N \chemfragphantom { \exp_not:n {#1} } } }
%% Protected, so the generated code survives any further full expansion.
\cs_new_protected:Npn \chemfragphantom #1 { \phantom {#1} }

%% Substituent on atom <n>, as a chemfig branch "( ... )".
%%   #1 = key list, #2 = atom number, #3 = relative angle (empty = let chemfig
%%   choose, as it does inside rings).
\cs_new:Npn \__chemfrag_branch:nnn #1#2#3
  { \exp_args:Ne \__chemfrag_branch_aux:nn { \__chemfrag_value:nnn {#1} {#2} { } } {#3} }
\cs_new:Npn \__chemfrag_branch_aux:nn #1#2
  { \tl_if_blank:nF {#1} { ( \__chemfrag_substituent:nn {#1} {#2} ) } }

%% A substituent may start with its own bond ("=O", "~N", "<:OH",
%% "-[::30]R").  Otherwise a single bond is added.  When an angle is
%% requested and the user did not give one, it is inserted after the bond.
\cs_new:Npn \__chemfrag_substituent:nn #1#2
  {
    \exp_args:Ne \__chemfrag_substituent_aux:nnn
      { \str_head:n {#1} } {#1} {#2}
  }
\cs_new:Npn \__chemfrag_substituent_aux:nnn #1#2#3
  {
    \__chemfrag_if_bond_char:nTF {#1}
      {
        \exp_args:Ne \str_if_eq:nnTF
          { \exp_args:Ne \str_head:n { \tl_tail:n {#2} } } { [ }
          { \exp_not:n {#2} }
          {
            \tl_to_str:n {#1}
            \tl_if_blank:nF {#3} { [ \c_colon_str \c_colon_str #3 ] }
            \tl_tail:n {#2}
          }
      }
      {
        -
        \tl_if_blank:nF {#3} { [ \c_colon_str \c_colon_str #3 ] }
        \exp_not:n {#2}
      }
  }
\prg_new_conditional:Npnn \__chemfrag_if_bond_char:n #1 { TF }
  {
    \bool_lazy_any:nTF
      {
        { \str_if_eq_p:nn {#1} { - } }
        { \str_if_eq_p:nn {#1} { = } }
        { \str_if_eq_p:ee {#1} { \c_tilde_str } }
        { \str_if_eq_p:nn {#1} { < } }
        { \str_if_eq_p:nn {#1} { > } }
      }
      { \prg_return_true: } { \prg_return_false: }
  }

%% Bond character <n> of a pattern such as "=-=-=-" ("-" when missing).
\cs_new:Npn \__chemfrag_bond_char:nn #1#2
  { \exp_args:Ne \__chemfrag_bond_char_aux:n { \tl_item:nn {#1} {#2} } }
\cs_new:Npn \__chemfrag_bond_char_aux:n #1
  { \tl_if_blank:nTF {#1} { - } { \tl_to_str:n {#1} } }

%% Expandable error for unknown fragments.
\msg_new:nnn { chemfrag } { unknown-ring }
  { Unknown~ring~`#1'.~Define~it~with~\iow_char:N\\chemfragdefinering. }
